\section{Introduction}

Generative Adversarial Networks (GANs) learn an adversarial game between a generator and discriminator and have become a standard foundation for high-fidelity synthesis \cite{GAN}. StyleGAN and StyleGAN2 further improved image quality by mapping an input latent code into an intermediate style space and modulating synthesis layers, making semantic manipulation more practical than in the original latent space \cite{karras2019stylebasedgeneratorarchitecturegenerative,karras2020analyzingimprovingimagequality}. These properties are useful for few-shot image generation, where a model must synthesize plausible samples for a novel category from very limited references rather than retraining a generator from scratch.

Latent editing methods exploit the observation that many attributes correspond to directions or subspaces in StyleGAN representations \cite{shen2020interpretinglatentspacegans,harkonen2020ganspacediscoveringinterpretablegan}. Attribute Group Editing (AGE) applies this idea to few-shot generation by separating category-relevant class structure from category-irrelevant variations such as color, pose, and texture \cite{AGE}. Stable Attribute Group Editing (SAGE) improves AGE by estimating a more stable class embedding from few-shot examples and by adaptively selecting compatible attribute directions from related seen classes \cite{ding2023stableattributegroupediting}. However, SAGE still depends on a pSp inversion path, and encoder-based inversion can trade reconstruction fidelity against editability \cite{richardson2021encodingstylestyleganencoder,liu2023delvingstyleganinversionimage}.

This study addresses that limitation by integrating the StyleFeatureEditor (SFE) Inverter and Feature Editor into SAGE. Instead of representing an input only through $W^{+}$ style codes, FeatureSAGE predicts both a $W^{+}$ latent and an intermediate feature tensor, then edits both during synthesis \cite{bobkov2024devildetailsstylefeatureeditordetailrich}. The research problem is whether this feature-level inversion path can improve detail preservation, edit stability, and distributional quality for novel classes while retaining the SAGE mechanism for class-consistent few-shot attribute group editing.

Specifically, the study asks whether replacing pSp with SFE improves inversion fidelity and image-detail preservation, whether the change affects consistency and diversity for unseen Flowers and Animal Faces classes, and how the proposed framework compares with AGE and SAGE in image quality and edit reliability. These questions are important because few-shot augmentation is useful only when generated samples preserve category identity while still providing realistic non-category variation. A generator that produces visually plausible but class-inconsistent images can reduce downstream value, while a model that preserves the reference too closely may fail to provide useful diversity.

The practical motivation is data-efficient image synthesis for fine-grained domains where collecting many labeled examples is difficult. Animal-face and flower categories contain meaningful intra-class variation in color, texture, orientation, and local structure, but they also require strong preservation of category-specific appearance. This makes them suitable testbeds for evaluating whether a few-shot generator can introduce useful variation without corrupting the class-defining features of the reference image.

The contributions of this paper are as follows: (1) a FeatureSAGE framework that replaces the pSp-based SAGE inversion path with a StyleFeatureEditor Inverter and Feature Editor; (2) a one-shot evaluation protocol on FUNIT Animal Faces and 102 Flowers using the same pretrained StyleGAN2 setting, data splits, and metrics across AGE, SAGE, and FeatureSAGE; and (3) quantitative and qualitative evidence showing that FeatureSAGE lowers FID while keeping LPIPS close to the baselines. The study is limited to StyleGAN2-based few-shot image generation on Animal Faces and Flowers and does not evaluate optimization-based inversion or non-StyleGAN generators.
